\section{Threats to Validity}
\label{sec:threats-to-validity}

In previous sections of this chapter, we address our research questions (RQ1--RQ3) 
(Section~\ref{subsec:eval-design-rqs})
with twelve targeted studies. The other three studies are used to evaluate the fairness, 
security and computational scalability of MERIT. Like any other evaluation, we need to consider that these
results are subject to threats that can detract from any of the claims that we make. We list these threats
according to the works of Wohlin et al.~\cite{wohlin_threats}. More specifically, we have the following forms of
validity to which we need to consider threats: construct, internal, external and conclusion validity.
In this section, we state what each metric can, and cannot, support. 

\subsection{Construct Validity}
\label{subsec:construct-validity}

The first form of validity is that of construct validity, which refers to the extent
to which the metrics and instruments we use actually measure the research questions we claim to address.

\textbf{Synthetic ground truth.}
Some of our studies, specifically Studies~01, and 07--10, use synthetic ground truth data in the form of
JSON profiles and CV PDFs. Due to the nature of this data being created by the author, it is subject to
confirmation bias, the evaluation data we synthesise may unintentionally be authored to exhibit the
behaviours that we claim to test for, essentially cherry-picking the data to support our claims.
We aim to reduce the effect of confirmation bias by deliberately including cases that are designed to
fail for our system to provide a more robust evaluation. One study worth pointing out is the Signal
Dissonance study (Study~10), where we purposely create contradictory multi-source evidence to 
contradict the CV. None of our synthetic data is externally validated, and we do not use
real-world data due to GDPR and PII constraints. The controlled personas and synthetic data
were, therefore, the most ethical way to stress-test our system without the risk of exposing real candidates,
while maintaining reproducibility, since we include the data in the GitHub repository.

\textbf{Proxy metrics and the gap between ingestion and scoring.}
Throughout our studies, we report metrics such as parser F$_1$, Spearman $\rho$, shortlist rates, SUS, and more.
These metrics are not hire outcomes, nor interview pass rates or offer acceptance rates. Because of this, we
cannot claim that MERIT ``improves hiring outcomes''. A high Spearman $\rho$ or narrow Blind Mode prestige gap
cannot be interpreted as MERIT improving these decisions, however, we can still observe the nuanced behaviour
of our system and draw specific conclusions about the behaviour of our system. Due to the small sample size,
we cannot generalise the results to a larger population.

Another thread is that our Human-Machine Alignment studies (Study~07--10) use pre-parsed JSON data that assumes
perfect ingestion of data. In our RQ1 studies, we show that ingestion is not perfect, even under authored synthetic data,
therefore end-to-end MERIT behavior may differ from what we observe in these studies as they do not account for parser noise.

\textbf{Study-specific construct threats (01, 04--05, 06--11, 13--15).}
Study~01 and 06--10 intentionally use sparse LinkedIn personas to stress test the
contrast between a CV's claimed evidence and GitHub's demonstrated evidence. LinkedIn as a data source is considered
supplementary to a CV profile to account for parser error. However, since we assume perfect ingestion of data, we do not
need LinkedIn to fill in the gaps, therefore we can still observe the behaviour of our system without the need for LinkedIn
(Section~\ref{subsec:study-01-comparative-ablation}, 
Section~\ref{sec:human-machine-alignment}).

Study~04's skill recall does not reflect parser failure alone, it reflects the dictionary scope of the Devicon dataset
(Section~\ref{subsec:study-04}).
Study~05 bounds how much we can trust JD-driven metric configuration when JD extraction is imperfect (Section~\ref{subsec:study-05}).
Study~13 employs a single global corpus for all TF-IDF calculations, meaning that the results suffer
from cross-role IDF interference. The study also suffers from uneven distribution of scores from the participants who
help in the generation of the ground truth. Since TF-IDF aims for a fairly even distribution of scores, we depress
the $\rho$ because of these decisions (Section~\ref{subsubsec:study-13-limitations}).

The SUS scores generated from Study~14 indicate prototype feasibility, they are 
not product-market validation (Section~\ref{subsec:hci-trial-protocol}).
Blind Mode standardises layout, redacts names, employers and institutions, but it keeps tier labels due to its
usage in the scoring pipeline. Because of this, Study~15 does not guarantee demographic fairness.
We do not mask engine scores in Blind Mode, so evaluators may have anchored on match percentages for shortlist decisions (Section~\ref{subsec:bias-mitigation-efficacy}).
Blind Mode must not be interpreted as a fairness certificate.

\subsection{Internal Validity}
\label{subsec:internal-validity}

Internal validity concerns whether the observed results are caused by the variables 
under study, rather than by confounding factors.

\textbf{In-repository baselines and benchmarks.}
The baselines in our studies (Traditional ATS and Modern AI ATS) are in-repository implementations that aim to
proxy the keyword-matching and AI-assisted ATS categories mentioned earlier in Section~\ref{sec:automated-screening}.
They are not commercial products, nor do they replicate the full complexity or sophistication of the real-world ATS systems.
The keyword baseline may produce lower scores than the semantic baseline in real-world ATS systems.
The Modern AI baseline uses SBERT to compute cosine similarities between CV embeddings and JD embeddings, real world
ATS systems may employ more sophisticated models which would likely produce better results.
Reported baseline statistics for runtime and spacetime compare the algorithm families on the author's workstation.
Our baselines do not account for factors such as API latency, PDF ingestion, and network I/O.
Studies~02--03 only record single runs on a workstation without confidence intervals (Section~\ref{subsec:studies-02-03-scalability}).
Scores are likely to vary from run to run, however, the general pattern will remain stable across 
configurations under identical hardware and software.


\textbf{Protocol and scenario design.}
As noted under construct validity, our engines use pre-parsed JSON for evaluation, here we note
that the rank shifts do not account for ingestion errors, they do not validate the full PDF upload pipeline.
Also, Study~14 always runs the first stage of the pipeline before the second, so the reversals and confidence
gains might be influenced by order effects, rather than just explainability.
Study~06 clones one honest baseline and perturbs it with eight different scenarios. This cleanly isolates
each perturbation, but it does not tell us how often such attacks appear in real candidate pools (Section~\ref{subsec:study-06-adversarial}).
Studies~01 and~10 intentionally include engineered success and failure narratives to aid explanation 
(Sections~\ref{subsec:study-01-comparative-ablation}, 
\ref{subsec:study-10-spearman-dissonance}).

\subsection{External Validity}
\label{subsec:external-validity}

This form of validity revolves around determining the extent to which these findings can be generalised to the wider population,
beyond the conditions of these studies.

\textbf{Sample size and controlled scope.}

The cohort sizes for Studies~07--10 are small, more specifically, 10 candidates per cohort. Due to this
both the $\rho$ and $p$-values are volatile, and we cannot draw any conclusions about the generalisability 
of the results based on the $p$-values.
Studies~13--15 involve small panels of $N = 4$--$5$ participants, thus it cannot support a generalisation
at the population level.
In our evaluations, we often prioritise controlled ablation over field realism
since we want to isolate the behaviour of our system from the noise of the real world.
Synthetic stress tests and in-repository baselines help us isolate MERIT's
inner mechanisms, however, they do not replicate live ATS workflows, nor real hiring outcomes.

\textbf{Templates, evidence, and deployment.}
Study~04's multi-column layouts follow a specific template family, they do not reflect the
full diversity of multi-column layouts in the real world. MERIT's multi-source pillar rewards
public LinkedIn and GitHub signals in the majority of cases, therefore, we cannot generalise 
the results to the full population of candidates. Our integration of LinkedIn is not deployment-ready,
we need to benchmark it against an official LinkedIn API to be able to generalise the results.
Candidates without public evidence (Carl Corp, Section~\ref{subsec:study-01-comparative-ablation})
are often disadvantaged relative to juniors that have visible open-source histories.
Benchmarking doesn't account for the latency introduced by API calls, PDF ingestion, Network I/O, and database persistence.
The measurements we collect are not representative of production latency and memory usage.

\subsection{Conclusion Validity}
\label{subsec:conclusion-validity}

Conclusion validity addresses whether the conclusions we draw from our data are actually justified
by our choice of method and metrics.

\textbf{Spearman $\rho$ as the primary alignment metric.}

Spearman $\rho$ is a measurement of monotonic agreement against another ordering. It cannot measure
the calibration of match percentages, nor can it measure the magnitude of individual rank swaps.
Our engines may produce similar $\rho$ values, but they may produce significantly different score spreads.
We include the full rankings and scores for each engine in the project repository~\cite{merit_repo}, rather than
treating correlation coefficients alone as enough evidence. Studies~07--10 have recruiters rank only CV PDFs, 
while MERIT (Full) has access to GitHub and LinkedIn data. Spearman $\rho$ therefore reflects protocol
asymmetry along with engine quality (Section~\ref{subsec:spearman-method}).
The negative $\rho$ in Study~08 is due to engine misalignment under an intern JD, not sampling noise.
Our engine does not account for, nor claim seniority-aware hiring unless we introduce explicit policy rules
to account for this.

\textbf{Shapley verification and multiple studies.}
In Study~11, we verify efficiency, null-player, additivity, and determinism on a set of ten candidate profiles.
We intentionally leave out the axiom of symmetry since sources are heterogeneous by design (Section~\ref{subsec:study-11-shapley-verification}).
The negative GitHub attribution for Vince Vault is mathematically valid in our evaluation, however it
may be considered as a cognitive surprise to recruiters.

We run fifteen studies on overlapping personas and engines, which creates a risk of selectively emphasising supportive
results, even if we document failures explicitly.
The chapter should not be read as a meta-analysis with pooled effect sizes.

\subsection{Bounded Claims}
\label{subsec:validity-mitigations}

Despite the limitations we note in this section, the chapter still provides convergent directional evidence under
controlled data.
We therefore frame three bounded claims, one per research question.
Each claim states what the evaluation supports under our internal criteria (Table~\ref{tab:eval-literature-traceability}), and what it does not.


\textbf{RQ1 (Extensibility).}
Studies~04--05 allow us to claim that PDF ingestion is usable, but below the $\geq 80\%$ parsing target.
Skill recall is bounded by the scope of the Devicon dataset (38.6\%). Other metrics from this study are
influenced by layout sensitivity, especially for names on multi-column CVs (Section~\ref{subsec:study-04}).
Study~13 shows that TF-IDF metric priorities align with recruiter judgement directionally (aggregate $\rho = 0.25$), 
presenting exploratory evidence only (Section~\ref{subsubsec:study-13-limitations}).


\textbf{RQ2 (Multi-source fusion).}
Study~12 confirms that Beta--Bernoulli fusion favours demonstrated evidence over self-reported claims (Section~\ref{subsec:study-12-bayesian-conflict-resolution}).
Studies~01 and~06 show rank and score shifts that match design intent for hidden talent, fraud and
adversarial personas (Sections~\ref{subsec:study-01-comparative-ablation} and~\ref{subsec:study-06-adversarial}).
Studies~07--10 show that MERIT (Full) meets the directional $\rho \geq 0.7$ criterion in Studies~07 and~09,
however, this is not the case in Studies~08 or~10, and it does not lead every cohort column (Section~\ref{subsec:spearman-synthesis}).
We claim that multi-source fusion and integrity checks can improve alignment with expert ordering when
supplementary evidence agrees with the CV. We do not claim that MERIT (Full) replaces keyword or semantic
baselines in all hiring contexts, nor that it can perform well when candidates have no public evidence.
We cannot claim that MERIT works well when presented with batch-wide contradictory evidence, nor can we
claim MERIT can perform well when presented with nuances such as overqualification (Section~\ref{subsec:study-08-spearman-seniority}).

\textbf{RQ3 (Explainability and Transparency).}
Study~11 verifies the mathematical correctness of the Shapley values, and shows 
that they satisfy the efficiency, null-player, additivity, and determinism axioms on synthetic profiles (Section~\ref{subsec:study-11-shapley-verification}).
Negative GitHub attributions are also mathematically valid in that formulation.
Study~14 meets the SUS target in our pilot ($N=5$, mean SUS 81.5), and we observe decision reversals
throughout the study, presenting directional evidence for glass-box feasibility (Section~\ref{subsec:hci-trial-protocol}).
Study~15 shows directional compression of prestige-driven shortlist gaps in the UI (50~pp to 12.5~pp), but it is not
proof of demographic fairness due to the small-scale pilot (Section~\ref{subsec:bias-mitigation-efficacy}).
We claim that our transparency makes penalties and uplifts contestable for recruiters, but we do not 
claim that it removes automation bias since visible engine scores and tier labels can still anchor 
evaluators rather than eliminate bias (Section~\ref{subsec:construct-validity}).

\textbf{Supplementary: Feasibility and Security.}
Studies~02--03 both support deployment feasibility from a scalability and performance perspective under
three data sources. We can claim that MERIT's Shapley explainability layer is usable, but it is not cost-free at scale.
It is scalable with candidate size, but not with the number of data sources.
Study~06 shows that adversarial perturbations can be used to stress test the integrity of the system, and that
the system is able to detect and respond to the majority of these perturbations. We do not claim that MERIT
defends against all adversarial attacks, nor can we claim that it is a fair hiring system.

\textbf{Final Claims.}
MERIT is not fairness-certified and should not be considered as a fully autonomous hiring authority.
Future work requires counterbalanced HCI trials to account for order effects.
To make stronger claims, MERIT needs larger independent recruiter panels, repeated benchmarks with variance reporting
on various hardware and language configurations. MERIT would also need end-to-end PDF ingestion tests as well as 
field studies that can link MERIT rankings to interview, hire and performance outcomes.

MERIT is best described as an \textbf{auditable screening assistant}, where its purpose is
to help recruiters inspect evidence. Decisions should not be treated as the final decision, but rather as a hypothesis
to be tested and challenged by the recruiter.
